\documentclass[11pt]{article}

% Page layout
\usepackage[letterpaper,left=1.5in,right=1in,top=1in,bottom=1in]{geometry}

% Font setup
\usepackage{iftex}
\ifPDFTeX
\usepackage[scaled]{helvet}
\renewcommand{\familydefault}{\sfdefault}
\else
\usepackage{fontspec}
\setmainfont{Arial}
\fi

% Core packages
\usepackage{amsmath}
\usepackage{graphicx}
\usepackage{multicol}
\usepackage{setspace}
\usepackage{tikz-cd}
\usepackage{fancyhdr}
\usepackage{booktabs}
\usepackage{array}
\usepackage{longtable}
\usepackage[numbers]{natbib}
\usepackage{hyperref}
\usepackage{qrcode}
\bibliographystyle{ACM-Reference-Format}

% Paragraph and spacing
\usepackage{indentfirst}
\setlength{\parindent}{0.5in}
\setlength{\parskip}{0pt}
\onehalfspacing

% Header/footer
\pagestyle{fancy}
\fancyhf{}
\fancyhead[R]{\thepage}
\renewcommand{\headrulewidth}{0pt}
\setlength{\headheight}{14pt}

\begin{document}

\begin{center}
  {\LARGE \textbf{Technical Report: Data Analysis on Titanic}}\\
  \vspace{0.5em}
  \textbf{Meinard Adrian P. Francisco}\\
  \textit{College of Computer and Information Sciences, Polytechnic University of the Philippines}\\
  \text{\today} % March 15, 2026
\end{center}

\section{Introduction}
The sinking of the RMS Titanic is one of the most infamous shipwrecks in history. On April 15, 1912, during her maiden voyage, the Titanic sank after colliding with an iceberg, killing 1502 out of 2224 passengers and crew. This tragedy shocked the international community and led to better safety regulations for ships. One of the reasons that the shipwreck led to such loss of life was that there were not enough lifeboats for the passengers and crew. Although there was some element of luck involved in surviving the sinking, some groups of people were more likely to survive than others, such as women, children, and the upper-class.

In this report, I objectively quantified and analyzed these claims using historical passenger data. The full implementation and source code for the Titanic analysis pipeline are available on GitHub at \url{https://github.com/znarfm/cttc-data-science}. To conduct a robust analysis problem, raw data undergoes a sequence of processing via a Python-based pipeline comprised of three modular scripts: initial data engineering (\texttt{clean\_data.py}), detailed cross-tabular analysis (\texttt{analyze\_data.py}), and feature charting (\texttt{visualize\_data.py}). These scripts isolate structural biases in survival likelihoods and visually emphasize them.

\section{Methodology}

\subsection{Data Source and Collection}
The primary dataset used is the publicly available Titanic passenger manifest from Kaggle (\url{https://www.kaggle.com/datasets/yasserh/titanic-dataset}), comprising a sample of 891 records. The raw dataset includes demographic and travel information, though initial records contained missing values, specifically in the \texttt{Age} and \texttt{Embarked} columns.

\subsection{Data Preprocessing}
A Python-based script (\texttt{clean\_data.py}) was applied to clean the data using the \texttt{pandas} library. Variables selected for their potential socio-demographic impact like \texttt{Sex} and \texttt{Embarked} were converted into numeric codes. Missing boarding ports were assigned to Southampton (\texttt{S}, the majority). Missing \texttt{Age} values were assigned using the median age of the passenger's specific gender and class (\texttt{Pclass}) group. Finally, a \texttt{FamilySize} feature was made by combining the number of siblings/spouses and parents/children aboard.

\subsection{Analytical and Statistical Methods}
Statistical summaries and visualizations were generated using \texttt{analyze\_data.py} and \texttt{visualize\_data.py} through \texttt{matplotlib} and \texttt{seaborn}. Grouped survival percentages across categorical variables (\texttt{Pclass}, \texttt{Sex}, \texttt{Embarked}, and \texttt{FamilySize}) were calculated to quantify and visually identify structural biases influencing survival likelihoods.

\section{Results}
The dataset includes 891 passengers with an overall survival rate of 38.38\%. On average, passengers were 29 years old and paid a fare of 32.20, with most traveling alone (median family size of 0). Figure~\ref{fig:distributions} provides a high-level summary of these population distributions.

\begin{figure}[htbp]
    \centering
    \includegraphics[width=0.85\textwidth]{feature_distributions.png}
    \caption{Distributions of survival, class, gender, embarkation, and age across the passenger dataset.}
    \label{fig:distributions}
\end{figure}

\subsection{Survival by Gender}
A distinct survival advantage existed for women passengers. The analytical breakdown shows that the female survival rate reached 74.20\%, whereas the male survival only had 18.89\%. Figure~\ref{fig:gender} illustrates that although males formed the majority, females are the huge portion of survivors.

\begin{figure}[htbp]
    \centering
    \includegraphics[width=0.45\textwidth]{survival_rate_gender.png}
    \caption{Comparison of survival rates between female and male passengers.}
    \label{fig:gender}
\end{figure}

\subsection{Survival by Socio-Economic Class}
Passenger class, acting as a proxy for socio-economic status, linearly shown survival probabilities. First-class passengers survived at a rate of 62.96\%, second-class at 47.28\%, while the third-class drops to a mere 24.24\%. Figure~\ref{fig:pclass} clearly shows how survival odds decreased as ticket class rank lowered.

\begin{figure}[htbp]
    \centering
    \includegraphics[width=0.45\textwidth]{survival_rate_pclass.png}
    \caption{Survival rate broken down by passenger class (1st, 2nd, and 3rd).}
    \label{fig:pclass}
\end{figure}

\subsection{Survival by Age Group}
Age demographics also displayed unequal survival rates. Figure~\ref{fig:age} presents a stacked histogram distinguishing survivors from casualties across age brackets. While adults in their 20s and 30s formed the majority of casualties, infants and young children were notably prioritized, reflecting a higher relative proportion of survivors within the youngest age bins.

\begin{figure}[htbp]
    \centering
    \includegraphics[width=0.45\textwidth]{survival_rate_age.png}
    \caption{Stacked histogram displaying survivors and casualties categorized by age brackets.}
    \label{fig:age}
\end{figure}

\subsection{Survival by Family Size}
Family size had a non-linear effect on survival. People traveling alone had a 30.35\% survival rate, while small families of 1 to 3 members had much better odds (up to 72.41\% for a family of 3). In contrast, larger families (4 or more members) saw much lower survival rates, dropping to 0\% for the largest families. Figure~\ref{fig:family} illustrates this trend.

\begin{figure}[htbp]
    \centering
    \includegraphics[width=0.45\textwidth]{survival_rate_family_size.png}
    \caption{Survival rate distributed across varying family sizes.}
    \label{fig:family}
\end{figure}

\section{Discussion}
The results confirm that survival on the Titanic was heavily influenced by social and economic factors rather than just luck. The high survival rate for women supports the historical \textit{women and children first} policy, while men faced much lower odds. First-class or higher class passengers likely had better access to lifeboats and earlier warnings, leading to much higher survival rates compared to those in lower or third class. Additionally, small families (1 to 3 members) were more successful, possibly because they could coordinate quickly without the logistical challenges of larger groups/families.

In conclusion, being female, traveling in first class, and having a small family were the greatest advantages during the disaster. These findings statistically reaffirm the historical record of events on how passengers were prioritized during the evacuation.

\end{document}
